\subsubsection{Roles of Financial Markets}

\begin{principle} \hlt{Core Market Functions}\\
A financial market is useful when it performs four linked functions:
\begin{enumerate}[label=\roman*.]
\setlength{\itemsep}{0pt}
\item \hlt{liquidity provision}: converts desired portfolio changes into executable trades;
\item \hlt{price discovery}: aggregates dispersed beliefs and information into observable prices;
\item \hlt{risk transfer}: reallocates exposures to agents more willing or able to bear them;
\item \hlt{capital allocation}: directs savings toward issuers and projects with attractive risk-adjusted uses of capital.
\end{enumerate}
Microstructure studies the frictions inside these functions rather than assuming they occur automatically.
\end{principle}

\begin{principle} \hlt{Immediacy Has a Price}\\
Asynchronous order arrival makes immediacy costly. The bid-ask spread compensates traders who absorb inventory, funding, processing, and adverse-selection risks \cite{demsetz_1968}.
\end{principle}

\begin{proposition} \hlt{Zero-Sum Trading Accounting}\\
For a trade at $P$ with common terminal value $V$, buyer and seller gains sum to zero: $(V-P)+(P-V)=0$. Frictions make trading negative-sum unless it creates information, liquidity, risk-transfer, or allocation value.
\begin{proof}
Buyer gain plus seller opportunity gain is zero before costs; costs subtract.
\end{proof}
\end{proposition}

\begin{definition} \hlt{Liquidity Provision}\\
Liquidity is the ability to trade size quickly, cheaply, and with limited price impact.\\
Supplied by limit orders, dealers, market makers, internalisers, and crossing systems, and consumed by marketable orders and urgent traders.
\end{definition}

\begin{mechanism} \hlt{Price Discovery}\\
Trades, quotes, cancellations, and depth transmit information about value, urgency, and risk-bearing capacity. Price discovery is strong when prices update quickly without being dominated by transient imbalances.
\end{mechanism}

\begin{mechanism} \hlt{Risk Transfer}\\
Markets separate cash-flow ownership from risk-bearing. Hedgers shed unwanted exposure; speculators, dealers, and market makers accept it for compensation.
\end{mechanism}

\begin{principle} \hlt{Capital Allocation}\\
Primary markets move capital from savers to issuers. Liquid and informative secondary markets provide exit prices, lower financing costs, and discipline poor uses of capital.
\end{principle}

\begin{principle} \hlt{Mechanism Relevance} \cite{o_hara_1995,madhavan_2000}\\
In Walrasian abstraction, the mechanism disappears after supply and demand clear. In real markets, sequencing, asymmetric information, costly immediacy, and priority rules make the mechanism part of price formation.
\end{principle}

\begin{caveat} \hlt{Microstructure Trade-Off}\\
Market design rarely improves every function at once. Transparency, spread size, venue opacity, and priority rules trade off liquidity provision, execution cost, information leakage, and public price discovery.
\end{caveat}
